\begin{abstractind}

Spektrum emisi merupakan cahaya yang terpancar dari elektron yang bertransisi antar tingkat energi atom atau ion. Struktur tingkat energi pada masing-masing atom atau ion membentuk pola unik yang memungkinkan analisis komposisi material secara presisi dalam metode \textit{Laser-Induced Breakdown Spectroscopy} (LIBS). Namun, akurasi metode konvensional ini seringkali terhambat oleh kompleksitas spektrum emisi multi-elemen, efek matriks, dan fluktuasi plasma. Penelitian ini mengembangkan pendekatan analisis berbasis \textit{Deep Learning} (DL) yang mempelajari spektrum emisi sintetis tersimulasi. Spektrum emisi sintetis disimulasikan menggunakan persamaan Saha dan distribusi Boltzmann, di mana setiap garis emisi dilebarkan menggunakan profil garis Gaussian dengan asumsi kondisi \textit{Local Thermal Equilibrium} (LTE) dan memanfaatkan data transisi atom dari \textit{Atomic Spectra Database} (ASD) NIST. Model Informer, sebuah arsitektur Transformer yang dioptimalkan dengan \textit{ProbSparse Self-Attention}, \textit{Self-Attention Distilling}, dan \textit{Generative-style Decoder}, dipilih untuk mengatasi sekuens spektrum emisi yang panjang. Model dilatih dan dievaluasi menggunakan metrik F1-Score untuk mencapai keseimbangan optimal antara presisi dan recall dalam identifikasi elemen. Pendekatan ini menyajikan bukti konsep yang kuat untuk deteksi multi-elemen yang andal dalam metode LIBS, serta berkontribusi bagi kemajuan spektroskopi modern.

\bigskip
\noindent
\textbf{Kata kunci:} LIBS, Informer, Spektrum Emisi Sintetis, Pembelajaran Mendalam, Analisis Prediktif, ProbSparse Self-Attention.
\end{abstractind}



